\section{Lattice design principles for a storage-ring EDM search}
\label{sec:principles}

\subsection{Spin motion in electromagnetic fields}

The concept of an EDM measurement in a storage ring is based on the
Thomas--Bargmann--Michel--Telegdi (T-BMT)
equation~\cite{Bargmann:1959gz,Fukuyama:2013ioa}. In accordance with the
Ehrenfest theorem, it describes the classical spin motion of a charged particle
with a hypothetical EDM:
\begin{equation}
\begin{aligned}
  \frac{d\vec{S}}{dt} &= \left(\vec{\Omega}_{mdm} + \vec{\Omega}_{edm}\right)\times\vec{S},\\
  \vec{\Omega}_{mdm} &= -\frac{e}{m\gamma}\left\{(\gamma G+1)\vec{B}_{\perp}
     + (1+G)\vec{B}_{\parallel}
     - \left(\gamma G+\frac{\gamma}{\gamma+1}\right)\frac{\vec{\beta}\times\vec{E}}{c}\right\},\\
  \vec{\Omega}_{edm} &= -\frac{e\eta}{2m}\left(\vec{\beta}\times\vec{B}+\frac{\vec{E}}{c}\right),
  \qquad G = \frac{g-2}{2},
\end{aligned}
\label{eq:tbmt}
\end{equation}
where $\gamma$ is the Lorentz factor, $\vec{\beta}$ is the relative velocity,
$c$ is the speed of light, $e$ and $m$ are the particle charge and mass, $G$ is
the magnetic anomaly, $g$ is the gyromagnetic ratio, $\Omega_{mdm}$ and
$\Omega_{edm}$ are the spin precession frequencies due to the magnetic and
electric dipole moments, $\eta$ is a dimensionless coefficient defined by
$d = \eta e\hbar/4mc$, and $\vec{B}=\{B_\perp, B_\parallel\}$ and
$\vec{E}=\{E_\perp, E_\parallel\}$ are the magnetic and electric fields. Since
no elements with a longitudinal magnetic field are used, we set
$B_\parallel = 0$; the longitudinal electric field is neglected as well,
$E_\parallel = 0$, because its contribution is small.

The lattice must provide conditions under which the EDM signal continuously
accumulates while the beam circulates in the ring. To this end, the spin and
the longitudinal momentum must remain aligned, $\vec{S}\uparrow\uparrow\vec{p}$,
with a precision that determines the attainable sensitivity of the EDM
measurement. Depending on the precision and the method of its realization,
this condition is referred to as the frozen spin or, more generally, the
quasi-frozen spin condition~\cite{Senichev:2016rez}.

\subsection{Frozen spin concept}

The idea of searching for the proton and deuteron EDMs with polarized beams in
a storage ring is based on the frozen spin (FS) method originally proposed at
the Brookhaven National Laboratory~\cite{Farley:2004hp,srEDM:2008}. An FS
lattice consists of deflectors combining electric and magnetic fields in one
element, in which the spin of the reference particle always stays along the
momentum.

The frozen spin method relies on the fact that at a specific, so-called
``magic'' energy the spin precesses in the external fields with the same
frequency as the particle momentum,
$\Omega^{p}_{E,B} = \frac{eE}{m\gamma\beta c} + \frac{eB}{m\gamma}$.
Subtracting this frequency from $\vec{\Omega}_{mdm}$ gives the spin precession
frequency relative to the momentum:
\begin{equation}
  \vec{\Omega}^{p}_{mdm} = \vec{\omega}^{p}_{E} + \vec{\omega}^{p}_{B},
  \label{eq:omega_p}
\end{equation}
where
$\vec{\omega}^{p}_{E} = \frac{e}{m}\left(G-\frac{1}{\gamma^2-1}\right)\frac{\vec{E}\times\vec{\beta}}{c}$
and $\vec{\omega}^{p}_{B} = \frac{e}{m}G\,\vec{B}_{\perp}$ are the
contributions of the electric and magnetic fields. In a purely electrostatic
ring ($\vec{B}=0$) the frozen spin condition $\vec{\omega}^{p}_{E} = 0$ is
satisfied at the magic energy
\begin{equation}
  G - \frac{1}{\gamma_{mag}^2-1} = 0 .
  \label{eq:magic}
\end{equation}
In a ring with both magnetic and electric elements, the frozen spin condition
$\vec{\Omega}^{p}_{mdm}=0$ is achieved by balancing the radial electric field
$E_r$ against the vertical guiding magnetic field $B_v$~\cite{srEDM:2008}:
\begin{equation}
  E_r = \frac{G B_v c\beta\gamma^2}{1-G\beta^2\gamma^2} \approx G B_v c\beta\gamma^2 .
  \label{eq:fs_hybrid}
\end{equation}
A hybrid ring is required for deuterons because their magnetic anomaly is
negative, $G = -0.143$, and condition~\eqref{eq:magic} cannot be satisfied.

It is convenient to introduce the spin tune, i.e., the number of spin
precessions per revolution. In an electrostatic ring the spin tune relative
to the momentum is $\nu_s^E = \omega^p_E/\Omega^p_E$:
\begin{equation}
  \nu_s^E = \left(G-\frac{1}{\gamma^2-1}\right)\gamma\beta^2 ,
  \label{eq:nu_e}
\end{equation}
and in a magnetic field, $\nu_s^B = \omega^p_B/\Omega^p_B$:
\begin{equation}
  \nu_s^B = \gamma G .
  \label{eq:nu_b}
\end{equation}
Thus, the same concept applies to protons and deuterons, but it is realized
with different types of deflectors.

\subsection{Quasi-frozen spin concept}

In the NICA ring, an implementation of the FS concept would require a
complete rebuild of the optics. Suppose instead that the spin oscillates in
the horizontal plane around the frozen spin direction with a small amplitude,
$\Phi_s^2 \ll 1$. This can be done with special electric-field elements in the
straight sections, which bring the spin back along the momentum after it has
deviated in the magnetic arc. The EDM signal growth is then reduced by the
factor $J_0(\Phi_s)\approx 1-\Phi_s^2/4$. The deuteron magnetic anomaly
$G=-0.143$ is small, and the spin oscillates around the momentum within half of
the spin phase advance in the magnetic arc,
$\Phi_s = \pi\gamma G/2 \approx 0.25$. Therefore, at the optimal parameters of
the EDM measurement, $\gamma = 1.12$, the effective EDM signal is reduced by
only a few percent.

This gives the essence of the quasi-frozen spin
concept~\cite{Senichev:2016rez}: the spin is not frozen relative to the
momentum but continuously oscillates around an average fixed direction that
coincides with the momentum direction. The question to be answered next is how
to implement a variable MDM spin precession in the storage ring, which is
needed for its zero average over the statistics, while providing a sufficient
growth of the EDM signal.

Two QFS options exist. In the first one, the magnetic and electric fields are
completely separated in space: the magnetic field acts in the arcs and the
electric field acts in the straight sections, which are realized as
negative-curvature ``reverse'' electrostatic arcs. This option inherits the
drawback of cylindrical electrodes, namely, a whole set of high-order field
nonlinearities. In the second option, a small magnetic field of about
\SI{100}{mT} is added to the spin-aligning E+B elements; it compensates the
Lorentz force of the electric field and makes the elements straight. Such an
E+B element is a Wien filter. The second option is the most appropriate one
for rings not designed specifically for the EDM search, in particular for the NICA ring
with bypass sections, whereas for the Nuclotron, where the free
space is limited, the first option is the baseline.

\subsection{Parameters of the Wien filters}
\label{sec:wien}

In a magnetic arc the momentum is rotated by the angle $\Phi^B_{arc}$, and the
spin simultaneously rotates in the horizontal plane relative to the momentum
by the angle
\begin{equation}
  \Phi_s^{arc} = \gamma G\,\Phi^B_{arc} .
  \label{eq:phi_arc}
\end{equation}
In the E+B elements of the straight section, the spin rotation relative to the
momentum has two components. In the electric field it is opposite to the
momentum rotation,
\begin{equation}
  \Phi_s^{E} = -\left(\gamma G+\frac{\gamma}{\gamma+1}\right)\beta^2\,\Phi^{E}_{ss},
  \label{eq:phi_e}
\end{equation}
where $\Phi^E_{ss}$ is the momentum rotation angle in the electric field. In
the magnetic field the spin rotates in the same direction as the momentum,
$\Phi_s^B = (\gamma G+1)\,\Phi^B_{ss}$, where $\Phi^B_{ss}$ is the momentum
rotation angle in the magnetic field. Since the Lorentz force vanishes, the
two momentum rotation angles are equal, $\Phi^E_{ss} = \Phi^B_{ss}$, and both
spin rotations can be expressed through one of them, e.g., through
$\Phi^B_{ss} = \frac{eB_{ss}}{m\gamma v}L_{ss}$, where $B_{ss}$ and $L_{ss}$
are the magnetic field and the length of the straight element. The QFS
condition requires that the net spin rotation in the E+B elements compensate
the spin rotation in the arc, $\Phi_s^B + \Phi_s^E = \Phi_s^{arc}$. For an arc
bending the beam by $\Phi^B_{arc}=\pi$ this gives
\begin{equation}
  (\gamma G+1)\,\Phi^B_{ss} - \left(\gamma G+\frac{\gamma}{\gamma+1}\right)\beta^2\,\Phi^E_{ss} = \pi\gamma G .
  \label{eq:qfs}
\end{equation}
Simple transformations yield the parameters of the E+B elements:
\begin{equation}
  L_\Sigma E_{ss} = \frac{G}{G+1}\,\frac{mc^2}{e}\,\pi\beta^2\gamma^3,
  \qquad
  B_{ss} = -\frac{E_{ss}}{c\beta},
  \label{eq:wien}
\end{equation}
where $L_\Sigma$ is the total length of the straight elements in one straight
section. For a realistic electric field of \SI{100}{kV/cm}, the corresponding
magnetic field is below \SI{96}{mT}. This simplifies the technical design: a
permanent magnet or an air-core coil can be used.

The material of this section has been reported
in~\cite{Senichev:2016rez,Senichev:2024qtr,Melnikov:2025gdj,Melnikov:2025ytd,Senichev:2025jqf}.
